\section*{الفصل \ref{ch:g11:organic-skeletons} --- الجزيئات العضوية: الهياكل والأسماء}

\begin{solution}{exo:g11:organic-skeletons:1}
في كل منهما ست \omterm{def:g7:atoms-and-molecules:atom}{ذرات} كربون (أطراف وزوايا) والصيغة \ce{C6H14}: فالهكسان
و2-ميثيل بنتان متماكبان.
\end{solution}

\begin{solution}{exo:g11:organic-skeletons:2}
البروبان؛ الهكسان؛ 2-ميثيل بيوتان.
\end{solution}

\begin{solution}{exo:g11:organic-skeletons:3}
\ce{CH3-CH2-CH3}؛ \chemfig{CH_3-CH(-[2]CH_3)-CH_3}؛
\chemfig{H-C(-[2]H)(-[6]H)-C(-[2]H)(-[6]H)-H}.
\end{solution}

\begin{solution}{exo:g11:organic-skeletons:4}
الأول، 3-ميثيل بنتان: 6 \omterm{def:g7:atoms-and-molecules:atom}{ذرات} كربون؛ وأطرافه \ce{CH3} (وهي ثلاثة)
تحمل 9 \omterm{def:g7:atoms-and-molecules:atom}{ذرات} هيدروجين، وزاويتاه \ce{CH2} تحملان 4، ونقطة التفرع
\ce{CH} تحمل 1: \ce{C6H14}. والهكسان الحلقي: 6 زوايا، كل منها
\ce{CH2}: \ce{C6H12}.
\end{solution}

\begin{solution}{exo:g11:organic-skeletons:5}
\ce{C8H18}. و$2n + 2 = 20$ يعطي $n = 9$: النونان، \ce{C9H20}.
\end{solution}

\begin{solution}{exo:g11:organic-skeletons:6}
\setchemfig{atom sep=1.6em}
3-ميثيل هكسان \chemfig{-[:30]-[:-30]-[:30](-[:90])-[:-30]-[:30]-[:-30]}،
\ce{C7H16}؛ 2,2-ثنائي ميثيل بيوتان \chemfig{-[:0](-[:90])(-[:-90])-[:0]-[:30]}،
\ce{C6H14}؛ 3-إيثيل بنتان
\chemfig{-[:30]-[:-30](-[:-90]-[:-30])-[:30]-[:-30]}، \ce{C7H16}.
\end{solution}

\begin{solution}{exo:g11:organic-skeletons:7}
البنتان و2-ميثيل بيوتان و2,2-ثنائي ميثيل بروبان، وهي مرسومة في شكل
الفصل.
\end{solution}

\begin{solution}{exo:g11:organic-skeletons:8}
«2-إيثيل بيوتان» له أطول سلسلة من خمس \omterm{def:g7:atoms-and-molecules:atom}{ذرات} كربون تمر بمجموعة
الإيثيل: إنه 3-ميثيل بنتان. و«4-ميثيل بنتان» مرقّم من الطرف
الخطأ: إنه 2-ميثيل بنتان. و«1-ميثيل بيوتان» ليس إلا سلسلة من خمس
\omterm{def:g7:atoms-and-molecules:atom}{ذرات} كربون: إنه البنتان.
\end{solution}

\begin{solution}{exo:g11:organic-skeletons:9}
2,2-ثنائي ميثيل بروبان (\qty{9.5}{\celsius}) < 2-ميثيل بيوتان
(\qty{27.8}{\celsius}) < البنتان (\qty{36.1}{\celsius}). \omterm{def:g7:atoms-and-molecules:atom}{فالذرات} نفسها
مرتبة على نحو أكثر فأكثر تراصًّا: تلامس أقل بين \omterm{def:g7:atoms-and-molecules:molecule}{الجزيئات}،
وتآثرات فان دير فالس أضعف، ودرجة غليان أدنى.
\end{solution}

\begin{solution}{exo:g11:organic-skeletons:10}
إغلاق السلسلة في حلقة يصنع رابطة \ce{C-C} إضافية، تحل
محل رابطتين \ce{C-H} (واحدة عند كل طرف). فالصيغتان الجزيئيتان مختلفتان:
ليسا متماكبين. والهكسان الحلقي ليس ألكانًا (فهيكله
حلقة؛ وصيغته ليست \ce{C_{n}H_{2n+2}}).
\end{solution}

\begin{solution}{exo:g11:organic-skeletons:11}
البيوتان و2-ميثيل بروبان (\ce{C4H10}). أما البروبان فصيغته \ce{C3H8} والبيوتان
الحلقي \ce{C4H8}.
\end{solution}

\begin{solution}{exo:g11:organic-skeletons:12}
\setchemfig{atom sep=1.5em}
\begin{center}
\small
\begin{tabular}{ccc}
\chemfig{-[:30]-[:-30]-[:30]-[:-30]-[:30]} &
\chemfig{-[:30](-[:90])-[:-30]-[:30]-[:-30]} &
\chemfig{-[:30]-[:-30](-[:-90])-[:30]-[:-30]} \\
الهكسان & 2-ميثيل بنتان & 3-ميثيل بنتان \\[6pt]
\chemfig{-[:0](-[:90])(-[:-90])-[:0]-[:30]} &
\chemfig{-[:30](-[:90])-[:-30](-[:-90])-[:30]} & \\
2,2-ثنائي ميثيل بيوتان & 2,3-ثنائي ميثيل بيوتان &
\end{tabular}
\end{center}
\end{solution}

\begin{solution}{exo:g11:organic-skeletons:13}
يستطيع خطان متعرجان طويلان أن يتجاورا على طولهما كله، فتتلامس \omterm{def:g7:atoms-and-molecules:atom}{ذرات}
كثيرة؛ أما الكرتان فلا تتلامسان إلا على مساحة صغيرة. وتآثرات فان دير
فالس تعمل بين \omterm{def:g7:atoms-and-molecules:atom}{الذرات} المتلامسة: فهي أكبر بين \omterm{def:g7:atoms-and-molecules:molecule}{الجزيئات} الخطية، التي تحتاج
إلى درجة حرارة أعلى لتنفصل.
\end{solution}

\begin{solution}{exo:g11:organic-skeletons:14}
أطول سلسلة فيها ست \omterm{def:g7:atoms-and-molecules:atom}{ذرات} كربون؛ وبالترقيم من اليسار تكون المستبدِلات
عند 2، 4، 4، ومن اليمين عند 3، 3، 5؛ والفرق الأول (2 مقابل 3) هو
الحاسم: 2,4,4-ثلاثي ميثيل هكسان.
\end{solution}

\begin{solution}{exo:g11:organic-skeletons:15}
\setchemfig{atom sep=1.6em}
\chemfig{-[:0](-[:90])(-[:-90])-[:-30]-[:30](-[:90])-[:-30]}: سلسلة من
خمس \omterm{def:g7:atoms-and-molecules:atom}{ذرات} كربون عليها مجموعتا ميثيل على الكربون 2 وواحدة على الكربون 4،
أي ثماني \omterm{def:g7:atoms-and-molecules:atom}{ذرات} كربون في المجموع؛ و$2 \times 8 + 2 = 18$ \omterm{def:g7:atoms-and-molecules:atom}{ذرة} هيدروجين،
\ce{C8H18}، وهي صيغة الأوكتان: فهما متماكبان. \omterm{def:g11:organic-skeletons:formulas}{والصيغة نصف
البنيوية}: \chemfig{CH_3-C(-[2]CH_3)(-[6]CH_3)-CH_2-CH(-[2]CH_3)-CH_3}.
\end{solution}

\begin{solution}{pb:g11:organic-skeletons:1}
\setchemfig{atom sep=1.5em}
\textbf{1.} \ce{C4H10}: مع $n = 4$، يكون $2n + 2 = 10$.

\textbf{2.} \chemfig{H-C(-[2]H)(-[6]H)-C(-[2]H)(-[6]H)-C(-[2]H)(-[6]H)-C(-[2]H)(-[6]H)-H}.

\textbf{3.} \ce{CH3-CH2-CH2-CH3}؛ \chemfig{-[:30]-[:-30]-[:30]}.

\textbf{4.} \chemfig{H-C(-[2]H)(-[6]H)-C(-[6]H)(-[2,1.5]C(-[1]H)(-[3]H)(-[2]H))-C(-[2]H)(-[6]H)-H}؛
\chemfig{CH_3-CH(-[2]CH_3)-CH_3}؛ \chemfig{-[:30](-[:90])-[:-30]}. فله
\omterm{def:g7:atoms-and-molecules:atom}{الذرات} نفسها، أربع \omterm{def:g7:atoms-and-molecules:atom}{ذرات} كربون وعشر \omterm{def:g7:atoms-and-molecules:atom}{ذرات} هيدروجين، مرتبطة بترتيب
مختلف.

\textbf{5.} نعم: متماكبان بنيويان (هيكلان مختلفان).

\textbf{6.} الثلاثة كلها: فدرجات غليانها أدنى من
\qty{20}{\celsius}.

\textbf{7.} عند \qty{-5}{\celsius}، أي تحت درجة غليانه، يبقى البيوتان
سائلًا تحت الضغط العادي ولا يطلق إلا قليلًا من الغاز.

\textbf{8.} البروبان (\qty{-42}{\celsius}) و2-ميثيل بروبان
(\qty{-11.7}{\celsius}) يغليان تحت درجات حرارة الشتاء: فهما لا يزالان يعطيان
غازًا وفيرًا في البرد.

\textbf{9.} 2-ميثيل بروبان متفرع، أكثر تراصًّا: فتآثرات فان دير
فالس بين جزيئاته أضعف.

\textbf{10.} مع ثلاث \omterm{def:g7:atoms-and-molecules:atom}{ذرات} كربون يكون الهيكل الوحيد الممكن هو
السلسلة \ce{C-C-C}: وأي \omterm{def:g7:atoms-and-molecules:atom}{ذرة} كربون تُنقل إلى مكان آخر تعطي السلسلة نفسها.

\textbf{11.} 2-ميثيل بيوتان.

\textbf{12.} يجب أن تُرقَّم السلسلة من الطرف الأقرب إلى
المستبدِل: 2، لا 3.

\textbf{13.} 2، و3، و5.

\textbf{14.} الهبتان، \chemfig{-[:30]-[:-30]-[:30]-[:-30]-[:30]-[:-30]}.

\textbf{15.} 2-ميثيل هكسان \chemfig{-[:30](-[:90])-[:-30]-[:30]-[:-30]-[:30]}
و3-ميثيل هكسان \chemfig{-[:30]-[:-30](-[:-90])-[:30]-[:-30]-[:30]}.
و«4-ميثيل هكسان» هو 3-ميثيل هكسان مرقّمًا من الطرف الخطأ.

\textbf{16.} 2,2-ثنائي ميثيل بنتان، و2,3-ثنائي ميثيل بنتان،
و2,4-ثنائي ميثيل بنتان، و3,3-ثنائي ميثيل بنتان.

\textbf{17.} نعم، 3-إيثيل بنتان. ولو كانت مجموعة الإيثيل على الكربون 2 لصنعت
سلسلة من ست \omterm{def:g7:atoms-and-molecules:atom}{ذرات} كربون تمر بها: فيكون \omterm{def:g7:atoms-and-molecules:molecule}{الجزيء} 3-ميثيل هكسان.

\textbf{18.} 2,2,3-ثلاثي ميثيل بيوتان،
\chemfig{-[:0](-[:90])(-[:-90])-[:0](-[:90])-[:0]}.

\textbf{19.} $1 + 2 + 4 + 1 + 1 = 9$.

\textbf{20.} للهبتان 9 \omterm{def:g11:organic-skeletons:isomer}{متماكبات بنيوية}، بما فيها الهبتان نفسه.
\end{solution}
